\stfhIndex{I146}{Dandiankaigeluomofu-Weitengbubianliang}{单点开格罗莫夫--威腾不变量（one-point open Gromov--Witten invariant）}
\section{\termOnePointOpenGromovWittenInvariant{}}
\label{sec:one-point-open-gw-invariant}

单调拉格朗日子流形的一个良好性质，是不存在第 10.6 节所述的\emph{负多重覆叠问题}。
这是因为每个非常值伪全纯圆盘都具有正马斯洛夫指标。由此可以定义不依赖于相容殆复
结构选择的\emph{开\personGromov{}--\personWitten{}不变量}。

本节依照 (Oh95a) 的说明，构造这一\emph{单点不变量}并证明其辛不变性。

对相容殆复结构 $J_0$，定义
\[
 \widetilde{\mathcal M}(L;\beta;J_0)
 =\{w:(D^2,\partial D^2)\to(M,L)\mid
 \overline\partial_Jw=0,\ w(\partial D^2)\subset L,\ [w]=\beta\}
\]
及其商
\[
 \mathcal M(L;\beta;J_0)
 =\widetilde{\mathcal M}(L;\beta;J_0)/PSL(2,\mathbb R).
\]
回忆指标公式
\begin{equation}
 \operatorname{vir.dim}\widetilde{\mathcal M}(L;\beta;J_0)
 =n+\mu_L(\beta),
\label{eq:disc-moduli-virtual-dimension}
\end{equation}
其中 $\mu_L(\beta)$ 是丛对 $(w^*TM,(\partial w)^*L)$ 的马斯洛夫指标
$\mu(w^*TM,(\partial w)^*L)$，而 $w$ 是
$\widetilde{\mathcal M}(L;\beta;J_0)$ 的某个（因而任意一个）元素。
\footnote{译者注：原书两次把边界丛写作 $(\partial w)^*L$，而非通常的
$(\partial w)^*TL$；译文照录。}
我们还考察带一个标记点的 $J_0$-全纯圆盘模空间
$\mathcal M_1(L;\beta;J_0)$。

给定光滑道路 $J:[0,1]\to\mathcal J_\omega$，定义参数化模空间
\[
 \mathcal M^{\mathrm{para}}(L;\beta;J)
 =\coprod_{t\in[0,1]}\{t\}\times\mathcal M(L;\beta;J(t)),
\]
\[
 \mathcal M_1^{\mathrm{para}}(L;\beta;J)
 =\coprod_{t\in[0,1]}\{t\}\times\mathcal M_1(L;\beta;J(t)),
\]
把它们看作 $[0,1]$ 上的纤维化。下面给出紧性结果。

\begin{proposition}
\label{prop:minimal-maslov-two-disc-compactness}
设 $L\subset(M,\omega)$ 单调，其\termMinimalMaslovNumber{} $\Sigma_L\geq2$。
设 $\beta\in\pi_2(M,L)$ 满足 $\mu(\beta)=2$。则：
\begin{enumerate}[label=(\arabic*)]
\item 对任意相容殆复结构 $J_0$，$\mathcal M(L;\beta;J_0)$ 与
$\mathcal M_1(L;\beta;J_0)$ 都紧；
\item 对道路 $J:[0,1]\to\mathcal J_\omega$，参数化模空间
$\mathcal M^{\mathrm{para}}(L;\beta;J)$ 与
$\mathcal M_1^{\mathrm{para}}(L;\beta;J)$ 都紧。
\end{enumerate}
\end{proposition}

\begin{proof}
这是开情形稳定映射紧性定理的直接推论。这里只证明第一条。设
$w_i\in\widetilde{\mathcal M}(L;\beta;J_0)$ 是一个序列。需要证明 $w_i$ 有一个
子序列，在模去 $PSL(2,\mathbb R)$ 作用后收敛到稳定映射。

由于 $[w_i]=\beta$ 固定，
\[
 E_{J_0}(w_i)=\int w_i^*\omega=\omega(\beta),
\]
故序列 $w_i$ 具有一致能量上界。稳定映射紧性定理于是说明，可以选取一个仍记作
$w_i$ 的子序列，收敛到稳定映射
\[
 w_\infty=\sum_{\ell=1}^{N}v_\ell,
\]
其中每个 $v_\ell$ 或为 $J_0$-全纯圆盘，或为 $J_0$-全纯球面。若 $v_\ell$ 是球面，
以 $\mu(v_\ell)$ 表示 $2c_1(v_\ell)$；若它是边界落在 $L$ 上的圆盘，则以
$\mu(v_\ell)$ 表示其马斯洛夫指标。

由于 $\mu(\beta)=2$ 极小，且
\[
 2=\mu(w_\infty)=\sum_{\ell=1}^{N}\mu(v_\ell),
\]
单调性蕴含上述和中至多有一个非常值分支。另一方面，$w_\infty$ 的稳定性又蕴含其中
不能有常值分支。因此必有 $N=1$，从而 $w_i$ 经适当重参数化后确实收敛。

类似论证证明 $\mathcal M_1(L;\beta;J_0)$ 及参数化空间的紧性。
\end{proof}

下面回忆通过扰动殆复结构得到的横截性结果。设 $L\subset(M,\omega)$ 是任意紧单调
拉格朗日子流形，以
\[
 \widetilde{\mathcal M}^{\mathrm{inj}}(L;\beta;J_0)
 \subset\widetilde{\mathcal M}(L;\beta;J_0)
\]
表示由\termSomewhereInjective{}的 $J_0$-全纯圆盘组成的子集。

闭曲面情形的如下基本横截性结果已在第 10.4、10.5 节证明。

\begin{proposition}
\label{prop:disc-evaluation-transversality}
设 $x\in L$。存在稠密子集
$\mathcal J_\omega(L,x)\subset\mathcal J_\omega$，使对所有
$\beta\in\pi_2(M,L)$：
\begin{enumerate}[label=(\arabic*)]
\item $\widetilde{\mathcal M}^{\mathrm{inj}}(L;\beta;J_0)$ 横截；
\item $\operatorname{ev}_\beta:\mathcal M_1^{\mathrm{inj}}(L,\beta;J_0)\to L$
横截于 $x$。
\end{enumerate}
对 $J=\{J(t)\}_{0\leq t\leq1}$ 的参数化模空间，同样结论成立。
\end{proposition}

称这样的对子 $(J_0,x_0)$ 为\termTransversePair{}。对\termTransversePair{} $(J_0,x_0)$，命题
\ref{prop:minimal-maslov-two-disc-compactness} 与
\ref{prop:disc-evaluation-transversality} 蕴含：对任意满足 $\mu(\beta)=2$ 的
$\beta$，$\operatorname{ev}_\beta^{-1}(x_0)$ 是紧零维流形（若非空）。

\begin{lemma}
\label{lem:minimal-maslov-discs-somewhere-injective}
设 $J_0$ 是任意相容殆复结构，不必正则。设 $L\subset(M,\omega)$ 单调且
$\Sigma_L>0$，并取满足 $\mu(\beta)=\Sigma_L$ 的
$\beta\in\pi_2(M,L)$。则
\[
 \widetilde{\mathcal M}^{\mathrm{inj}}(L;\beta;J_0)
 =\widetilde{\mathcal M}(L;\beta;J_0),
\]
即对任意 $J_0\in\mathcal J_\omega$，
$\widetilde{\mathcal M}(L;\beta;J_0)$ 中所有 $J_0$-全纯圆盘都\termSomewhereInjective{}。特别地，命题\ref{prop:disc-evaluation-transversality} 对完整
模空间成立。
\end{lemma}

\begin{proof}
第一条来自\termMinimalMaslovNumber{} $\Sigma_L$ 的定义：由像的结构定理——球面情形的
定理 8.7.3，以及圆盘情形的 (Laz00)、(KO00)——无论 $J_0$ 正则与否，每个
$J_0$-全纯圆盘都必为单圆盘。第二条则由单调性假设与 $\mu(\beta)=\Sigma_L$ 立即
推出，因为它们排除了冒泡。
\end{proof}

记
\[
 \Phi_{J_0}(L;x_0)
 =\sum_{\beta:\,\mu(\beta)=2}\#\bigl(\operatorname{ev}_\beta^{-1}(x_0)\bigr).
\]
\emph{从此以后不再特别说明：对一般 $L$，计数取 $\mathbb Z_2$-系数；当拉格朗日子流形
$L$（相对）spin 时，使用 $\mathbb Z$-系数。}

\begin{lemma}
\label{lem:one-point-count-invariance}
设 $L\subset(M,\omega)$ 单调且 $\Sigma_L\geq2$。设 $(J_0,x_0)$ 与
$(J_1,x_1)$ 是\termTransversePair{}，并且 $[x_0]=[x_1]\in\pi_0(L)$。则
\[
 \Phi_{J_0}(L;x_0)=\Phi_{J_1}(L;x_1).
\]
\end{lemma}

\begin{proof}
固定嵌入道路 $\gamma:[0,1]\to L$，其中 $\gamma(0)=x_0$、$\gamma(1)=x_1$。
由评价横截性的参数化版本及横截性延拓定理，可选道路
$J:[0,1]\to\mathcal J_\omega$，使得对所有 $\mu(\beta)=2$ 的 $\beta$，
$\widetilde{\mathcal M}^{\mathrm{para}}(L;\beta;J)$ 横截，且评价映射
\[
 \begin{aligned}
 \operatorname{Ev}_\beta:
 \mathcal M_1^{\mathrm{para}}(L;\beta;J)&\longrightarrow[0,1]\times L,\\
 (t,(w,z))&\longmapsto(t,w(z))
 \end{aligned}
\]
横截于 $\Lambda=\operatorname{id}\times\gamma:[0,1]\to[0,1]\times L$。指标公式给出
\[
 \dim\mathcal M_1^{\mathrm{para}}(L;\beta;J)=n+1,
\]
故 $\operatorname{Ev}_\beta^{-1}(\Lambda)$ 是紧一维流形，其边界为
\[
 \{0\}\times\operatorname{ev}_\beta^{-1}(x_0)
 \coprod
 \{1\}\times\operatorname{ev}_\beta^{-1}(x_1).
\]
由紧一维流形的分类定理，
\[
 \#\bigl(\operatorname{ev}_\beta^{-1}(x_0)\bigr)
 =\#\bigl(\operatorname{ev}_\beta^{-1}(x_1)\bigr).
\]
依定义即得结论。
\end{proof}

由此可以如下定义单点 GW 不变量。

\begin{definition}
\label{def:one-point-open-gw-invariant}
设 $L\subset(M,\omega)$ 是紧单调拉格朗日子流形，$\Sigma_L\geq2$；以
$\pi\in\pi_0(L)$ 表示 $L$ 的一个连通分支。相对于 $\pi$ 的单点 GW 不变量定义为
\[
 \Phi(L;\pi)=\Phi_{J_0}(L;x_0),
\]
其中 $(J_0,x_0)$ 是满足 $[x_0]=\pi\in\pi_0(L)$ 的\termTransversePair{}。当 $L$ 连通时，简记
\[
 \Phi(L;\{\mathrm{pt}\})=\Phi_{J_0}(L;x_0),
\]
称其为 $L\subset(M,\omega)$ 的单点 GW 不变量。
\end{definition}

下面考察几个例子。

\begin{example}
\label{ex:one-point-open-gw-examples}
\begin{enumerate}[label=(\arabic*)]
\item \stfhPropertyLabel{（标准环面）} 考虑标准环面
\[
 T^n=\underbrace{S^1(1)\times\cdots\times S^1(1)}_{n\text{ 个}}
 \subset\mathbb C^n\cong\mathbb R^{2n}.
\]
有 $\pi_2(\mathbb C^n,T^n)\cong\pi_1(T^n)\cong\mathbb Z^n$。以
$\beta_i$（$i=1,\ldots,n$）表示由映射
\[
 \begin{aligned}
 w_i:(D^2,\partial D^2)&\longrightarrow(\mathbb C^n,T^n),\\
 w_i(z)&=(1,\ldots,\underbrace{z}_{i\text{th}},\ldots,1)
 \end{aligned}
\]
代表的 $\pi_2(\mathbb C^n,T^n)$ 的显然生成元。计算其马斯洛夫指标与辛面积可知，
该环面单调，\termMonotonicityConstant{}为 $\lambda=\pi$。这些映射显然是经过点
$x_0=(1,\ldots,1)$ 的全纯映射，故至少有 $n$ 个彼此不同、马斯洛夫指标为 $2$
的全纯圆盘。

\item \stfhPropertyLabel{（\personClifford{}环面）} 这次考虑\personClifford{}环面
$T^n\subset\mathbb{CP}^n$。它由 $\mathbb C^{n+1}$ 上乘积圆周作用的商
$T^{n+1}/S^1\subset S^{2n+1}/S^1\cong\mathbb{CP}^n$ 定义。在射影坐标中有全纯
圆盘
\[
 w_j:(D^2,\partial D^2)\to(\mathbb{CP}^n,T^n),
 \qquad z\longmapsto[1;\ldots;\underbrace{z}_{j\text{th}};\ldots;1]
\]
（$j=1,\ldots,n+1$）。把
$\Phi(T^n,\mathbb{CP}^n;\{\mathrm{pt}\})=n+1$ 的证明留作习题。

\item \stfhPropertyLabel{（\personChekanov{}环面）} 对标准圆周
$S^1(1)\subset\mathbb R^2\cong\mathbb C$ 施行\personChekanov{}悬挂，得到\personChekanov{}环面
\[
 \Theta_0(S^1(1))\subset\mathbb R^2,
 \qquad
 \Theta_0(S^1(1))=I_1(N_0\times S^1(1)).
\]
\footnote{译者注：原书把所得\personChekanov{}环面的环境写作 $\mathbb R^2$；但紧接着定义的
辛嵌入以 $T^*\mathbb R^2$ 为靶，维数提示这里可能有环境维数的排印问题。译文照录。}
这里 $N_0$ 是 $T^*S^1$ 的零截面，$I_1=(i_1^*)^{-1}$ 是辛嵌入
\[
 I_1=(i_1^*)^{-1}:T^*(\mathbb R\times S^1)\longrightarrow T^*\mathbb R^2,
\]
其中 $i_1:S^1\times\mathbb R\to\mathbb R^2$ 定义为
\[
 i_1(\theta,q)=(e^q\cos(2\pi\theta),e^q\sin(2\pi\theta)).
\]
\personEliashberg{}与\personPolterovich{}在 (EP97) 中证明
\[
 \Phi(\Theta_0(S^1(1));\{\mathrm{pt}\})=1.
\]
不久前，\personChekanov{}与\personSchlenk{}尝试对若干显式辛流形类上的单调拉格朗日子流形集合
作完整分类 (CS10)。
\end{enumerate}
\end{example}

\begin{exercise}
\label{ex:classify-minimal-maslov-discs-tori}
证明例\ref{ex:one-point-open-gw-examples} 中给出的 $\mathbb C^n$ 内标准环面及
$\mathbb{CP}^n$ 内\personClifford{}环面的圆盘全都横截，并且它们给出了经过环面中给定点
$x_0$ 的所有马斯洛夫指标为 $2$ 的全纯圆盘。（若计入这些圆盘的定向，则它们全都
正贡献，故
$\Phi(T^n,\mathbb C^n;\mathrm{pt})=n$ 且
$\Phi(T^n,\mathbb{CP}^n;\{\mathrm{pt}\})=n+1$。定向研究见第 15.4 节。）
\end{exercise}

\stfhIndex{I054}{Qiekanuofu}{切卡诺夫（Chekanov）}
上述单点不变量计算立即导出\personChekanov{}定理 (Che96b)。

\begin{theorem}[\personChekanov{}]
\label{thm:standard-chekanov-tori-not-symplectically-equivalent}
标准环面 $T^2=S^1(1)\times S^1(1)$ 与\personChekanov{}环面
$\Theta_0(S^1(1))$ 不辛等价。
\end{theorem}

\personChekanov{}最初通过考察\personEkeland{}--\personHofer{}容量序列的行为证明该定理。应指出，标准环面与\personChekanov{}环面都满足下列条件：
\begin{enumerate}[label=(\arabic*)]
\item \termMinimalMaslovNumber{}为 $2$；
\item \termMonotonicityConstant{}为 $\pi$；
\item 第一\personEkeland{}--\personHofer{}容量为 $2\pi$。
\end{enumerate}
因此，没有简单的经典方法区分标准环面与\personChekanov{}环面。
